\documentclass[11pt,letterpaper]{article}

%% ── Packages ─────────────────────────────────────────────────────────────────
\usepackage[margin=1in]{geometry}
\usepackage{amsmath,amssymb,bm,bbm}
\usepackage{graphicx}
\usepackage{booktabs}
\usepackage{hyperref}
\usepackage{caption}
\usepackage{subcaption}
\usepackage{xcolor}
\usepackage{enumitem}
\usepackage{microtype}
\usepackage{parskip}
\usepackage{array}

\hypersetup{colorlinks=true, linkcolor=blue, citecolor=blue, urlcolor=blue}

%% ── Macros ───────────────────────────────────────────────────────────────────
\newcommand{\R}{\mathbb{R}}
\newcommand{\W}{\mathbf{W}}
\newcommand{\avec}{\mathbf{a}}
\newcommand{\psivec}{\boldsymbol{\psi}}
\newcommand{\norm}[1]{\left\|#1\right\|}

% Placeholder command for values to be filled after running experiments.
% Usage: \placeholder{description} renders as bold red [description].
\newcommand{\placeholder}[1]{\textbf{\textcolor{red}{[#1]}}}

\begin{document}

\noindent \textbf{ECE 66200: Pattern Recognition} \hfill \textbf{Spring 2026} \\
\vspace{1cm}

\begin{center}
    {\Large \textbf{Project 1, Part 3 (v2):\\
    Continual Deblurring via Filter Atom Swapping}} \\
    \vspace{0.5cm}
    {\large [Your Name]} \\
    \vspace{0.3cm}
    {\large \today}
\end{center}
\vspace{0.5cm}

%%=============================================================================
\section{Motivation}
\label{sec:motivation}
%%=============================================================================

When a neural network is trained sequentially on a stream of tasks, gradient
updates optimised for the current task overwrite weights that were critical for
earlier tasks --- a phenomenon known as \emph{catastrophic forgetting}
\cite{french1999catastrophic}.  The naive remedy is to maintain an independent
copy of the model for each task, but this incurs a memory cost that grows
linearly with the number of tasks and is impractical for deployment on
resource-constrained devices.

This problem is directly relevant to my research on \textbf{image deblurring
for spatially varying point spread functions (SV-PSFs)}.  A deployed imaging
system encounters many distinct blur conditions over its lifetime --- changes
in focus distance, aperture setting, atmospheric turbulence, or sensor aging
--- and each condition constitutes a new deblurring task.  A practical system
must master new blur conditions as they arise, without forgetting how to handle
previously encountered ones, and without storing a separate CNN for every blur
type.

The \emph{Decomposed Convolutional Filter} (DCF) framework
\cite{qiu2018dcfnet} provides a natural mechanism for this: by decomposing
each convolutional filter into shared \emph{atoms} $\psivec_k$ and
task-specific \emph{coefficients} $\avec_k$, we can isolate the task-dependent
and task-invariant components.  Building on the \emph{Filter Atom Swapping}
idea of Miao et al.\ \cite{miao2021continual}, we design a continual learning
system where only a handful of lightweight atom parameters are stored per task,
while the vast majority of the network --- coefficients, decoder, and batch
normalisation --- remains shared and frozen.

%%=============================================================================
\section{Problem Formulation}
\label{sec:formulation}
%%=============================================================================

\subsection{Split-MNIST Deblurring Tasks}

We construct five sequential deblurring tasks by partitioning the 10 MNIST
digit classes into pairs and degrading all images with the same Gaussian
point spread function.  Table~\ref{tab:tasks} defines the task sequence.

\begin{table}[h]
\centering
\caption{Five sequential Split-MNIST deblurring tasks.  All tasks share
         the same Gaussian blur ($\sigma = 1.5$, $7 \times 7$ kernel).}
\label{tab:tasks}
\begin{tabular}{clcc}
\toprule
Task $t$ & Digit classes & Training samples & Test samples \\
\midrule
0 & \{0, 1\} & $\sim$12{,}000 & $\sim$2{,}000 \\
1 & \{2, 3\} & $\sim$12{,}000 & $\sim$2{,}000 \\
2 & \{4, 5\} & $\sim$12{,}000 & $\sim$2{,}000 \\
3 & \{6, 7\} & $\sim$12{,}000 & $\sim$2{,}000 \\
4 & \{8, 9\} & $\sim$12{,}000 & $\sim$2{,}000 \\
\bottomrule
\end{tabular}
\end{table}

By keeping the blur kernel fixed across tasks, we isolate the
class-incremental learning challenge: the model must learn to reconstruct
new digit classes without forgetting previously learned ones, and the only
source of distributional shift between tasks is the spatial structure of the
digits themselves.  Figure~\ref{fig:blur_ex} shows example blurred and clean
images from each task.

\begin{figure}[h]
  \centering
  \includegraphics[width=0.85\textwidth]{fig_p3v2_blur_examples.pdf}
  \caption{Split-MNIST deblurring tasks.  \textbf{Top row:} Gaussian-blurred
    inputs ($\sigma = 1.5$).  \textbf{Bottom row:} clean targets.  Each column
    corresponds to one of the five tasks, with two digit classes per task.}
  \label{fig:blur_ex}
\end{figure}

\subsection{DCF Decomposition Recap}

Recall from Parts~1 and~2 that the DCF decomposition
\cite{qiu2018dcfnet} writes each convolutional filter as:
\begin{equation}
  \W[i,j] \;=\; \sum_{k=1}^{K} a_k[i,j]\,\psivec_k,
  \qquad
  \psivec_k \in \R^{l \times l},\quad a_k \in \R^{c' \times c},
  \label{eq:dcf}
\end{equation}
where $\psivec_k$ are \textbf{filter atoms} (a spatial basis dictionary) and
$a_k$ are \textbf{coefficients} (the task-specific mixing weights).  For
$K = 6$ atoms with $3 \times 3$ kernels, the atoms comprise only
$K \times l^2 = 6 \times 9 = 54$ parameters per convolutional layer.

\subsection{Forgetting Metric}

We measure catastrophic forgetting via \emph{backward transfer} (BWT):
\begin{equation}
  \mathrm{BWT} = \frac{1}{T{-}1}\sum_{j=0}^{T-2}
  \bigl(\mathrm{PSNR}_{T-1,\,j} - \mathrm{PSNR}^*_{j}\bigr),
  \label{eq:bwt}
\end{equation}
where $\mathrm{PSNR}^*_j$ is the peak PSNR on Task~$j$ measured immediately
after its own training, and $\mathrm{PSNR}_{T-1,\,j}$ is the PSNR on Task~$j$
after the final task has been trained.  Negative BWT indicates forgetting;
BWT $= 0$ means no performance degradation on any earlier task.

%%=============================================================================
\section{Design Choice: Filter Atom Swapping}
\label{sec:design}
%%=============================================================================

\subsection{Core Idea}

Our approach \emph{inverts} the coefficient-swapping strategy used in Part~2.
Rather than freezing atoms and swapping coefficients (which required storing
$\sim$93{,}000 parameters per task, yielding only 1.19$\times$ compression),
we \textbf{freeze the coefficients, decoder, and batch normalisation} after the
initial task and \textbf{swap only the atoms} for each new task.

The intuition is as follows.  The coefficients $\avec_k$ and the decoder
collectively define a general-purpose reconstruction pipeline: ``given a set of
spatial filter atoms, here is how to combine them into encoder filters and
how to decode the resulting feature maps back into a clean image.''  This
pipeline is learned on Task~0 and then kept fixed.  For each subsequent task,
only the atoms --- the spatial building blocks that the pipeline operates on
--- need to change to accommodate the new digit classes.

\subsection{Training Protocol}

\begin{enumerate}[leftmargin=*]
  \item \textbf{Task 0 (initial training):}
    Build a DCF autoencoder with $K = 6$ atoms per encoder layer
    (\texttt{bases\_grad=True}, Fourier-Bessel initialisation).
    Train \emph{all} parameters --- atoms $\psivec_k$, coefficients $\avec_k$,
    decoder weights, and batch normalisation --- end-to-end on Task~0 for
    30 epochs.  After convergence, store
    $\Psi^{(0)} = \{\psivec_k^{(0)}\}_{k=1}^{K}$ for all encoder layers.

  \item \textbf{Tasks $t = 1, 2, 3, 4$ (incremental):}
    \begin{itemize}
      \item \emph{Freeze:} coefficients $\avec_k$
            (\texttt{weight.requires\_grad = False}), all decoder parameters,
            and batch normalisation (set to eval mode so running statistics
            from Task~0 are preserved).
      \item \emph{Warm-start:} copy the Task~0 atoms $\Psi^{(0)}$ into the
            model.  This ensures the model starts from a coherent state
            (the coefficients were trained with these atoms).
      \item \emph{Train:} update \emph{only} the atoms $\psivec_k$
            on Task~$t$ for 30 epochs.
      \item \emph{Store:} save
            $\Psi^{(t)} = \{\psivec_k^{(t)}\}_{k=1}^{K}$ --- a snapshot
            of 162 scalars.
    \end{itemize}

  \item \textbf{Inference on Task $j$:}
    Load $\Psi^{(j)}$ into the encoder.  Because the coefficients, decoder,
    and batch normalisation are \emph{identical} to their Task-0 values, the
    resulting model is \emph{exactly} the one that was trained for Task~$j$.
    This guarantees \textbf{BWT $= 0.0$\,dB} by construction.
\end{enumerate}

\subsection{Memory Compression Analysis}

Let $P$ denote the total parameters of the DCF autoencoder and $P_\psi$ the
atom parameters across all encoder layers.  For our architecture:
\begin{align}
  P_\psi &= 3 \times (K \times l^2) = 3 \times 54 = 162, \\
  P_\text{shared} &= P - P_\psi.
\end{align}
For $T$ tasks:
\begin{equation}
  \text{DCF-CL memory} = P_\text{shared} + T \times P_\psi.
\end{equation}
The baseline (storing $T$ independent models of $P_\text{base}$ parameters
each) requires $T \times P_\text{base}$.  The compression ratio is therefore:
\begin{equation}
  \text{Compression} = \frac{T \times P_\text{base}}{P_\text{shared} + T \times P_\psi}.
  \label{eq:compression}
\end{equation}
Since $P_\psi \ll P$, the denominator is dominated by $P_\text{shared}$ and
the ratio grows linearly with $T$.  For $T = 5$ tasks in our experiment, the
expected compression is approximately
$\frac{5 \times P_\text{base}}{P_\text{shared} + 810}$, which evaluates to
roughly $5\times$.

\subsection{Why Atoms (Not Coefficients) Are Task-Specific Here}

This design choice may seem to contradict Part~2, where Exp~C showed that
\emph{coefficients} carry domain-specific information for MNIST$\to$SVHN
adaptation.  However, the two settings have fundamentally different task
structures:

\begin{itemize}[leftmargin=*]
  \item \textbf{Part 2 (domain adaptation):} Source and target share the same
    digit classes but differ in appearance (handwritten vs.\ street-view).
    The spatial patterns (edges, strokes) are similar across domains, so atoms
    are \emph{domain-invariant}.  The way these patterns are combined into
    class-discriminative features differs, making coefficients
    \emph{domain-specific}.

  \item \textbf{Part 3 (class-incremental reconstruction):} All tasks share
    the same image domain (MNIST) and the same blur, but the digit classes
    differ.  The reconstruction objective requires different spatial atoms to
    capture the distinct stroke geometries of each digit pair (e.g., ``0'' and
    ``1'' vs.\ ``8'' and ``9''), while the general reconstruction pipeline
    (coefficient combination + decoding) transfers across digit classes.
\end{itemize}

The appropriate swapping target therefore depends on the \emph{nature of the
task variation}: atoms for class-level differences, coefficients for
domain-level differences.

%%=============================================================================
\section{Experimental Settings and Results}
\label{sec:results}
%%=============================================================================

\subsection{Architecture and Training Details}

Both the DCF-CL and baseline models use the same autoencoder architecture
from Part~1 Task~3:

\begin{itemize}[leftmargin=*]
  \item \textbf{Encoder:} three convolutional layers
    ($1 \to 32 \to 64 \to 64$), each followed by BatchNorm and ReLU;
    MaxPool(2,2) after the first two layers.  The DCF-CL variant replaces
    Conv2d with Conv\_DCF ($K = 6$, \texttt{bases\_grad=True},
    Fourier-Bessel initialisation, mode1).
  \item \textbf{Decoder:} three ConvTranspose2d layers
    ($64 \to 64 \to 32 \to 1$) with BatchNorm, ReLU, and a final Sigmoid.
  \item \textbf{Training:} Adam optimiser, learning rate $10^{-3}$, cosine
    annealing, 30 epochs per task, batch size 128, MSE loss.
\end{itemize}

% ── NOTE TO AUTHOR: fill placeholder values from cl_deblur_results.json ──
% Search for \placeholder{...} and replace with actual numbers.

\subsection{PSNR Results}

Table~\ref{tab:final_psnr} reports the final PSNR on each task after all
five tasks have been trained.  The baseline suffers severe forgetting on
earlier tasks, while DCF-CL maintains its peak performance on every task
through exact atom restoration.

\begin{table}[h]
\centering
\caption{PSNR (dB) on each task after the final training phase (Task~4).}
\label{tab:final_psnr}
\begin{tabular}{lcccccl}
\toprule
Method & T0 $\{0,1\}$ & T1 $\{2,3\}$ & T2 $\{4,5\}$ & T3 $\{6,7\}$ & T4 $\{8,9\}$ & BWT (dB) \\
\midrule
Baseline   & \placeholder{B0} & \placeholder{B1} & \placeholder{B2} & \placeholder{B3} & \placeholder{B4} & \placeholder{BWT\_base} \\
DCF-CL     & \placeholder{D0} & \placeholder{D1} & \placeholder{D2} & \placeholder{D3} & \placeholder{D4} & \textbf{0.0} \\
\bottomrule
\end{tabular}
\end{table}

\subsection{Forgetting Curves}

Figure~\ref{fig:forget} displays the PSNR trajectory for each task across
the five training phases.  For the baseline (red, dashed), previously learned
tasks exhibit steep PSNR drops as new tasks are trained --- the hallmark of
catastrophic forgetting.  For DCF-CL (blue, solid), the PSNR for each
completed task remains perfectly flat across all subsequent training phases,
confirming zero backward transfer.

\begin{figure}[h]
  \centering
  \includegraphics[width=\textwidth]{fig_p3v2_forgetting_curves.pdf}
  \caption{Forgetting curves: PSNR on each task across training phases.
    \textbf{Baseline (red, dashed):} severe forgetting on earlier tasks as
    later tasks are trained.
    \textbf{DCF-CL (blue, solid):} perfectly flat curves for completed tasks,
    confirming BWT $= 0.0$\,dB.}
  \label{fig:forget}
\end{figure}

\subsection{PSNR Heatmaps}

Figure~\ref{fig:heatmap} presents the full $5 \times 5$ PSNR matrices.  In
the baseline heatmap (left), entries below the diagonal drop sharply,
quantifying the cumulative forgetting as more tasks are trained.  In the
DCF-CL heatmap (right), each column is constant from the task's training
phase onward, confirming that atom swapping provides an exact model restore
with no approximation error.

\begin{figure}[h]
  \centering
  \includegraphics[width=\textwidth]{fig_p3v2_psnr_heatmaps.pdf}
  \caption{PSNR matrices (\textbf{rows} = training phase;
    \textbf{cols} = evaluation task).
    \textbf{Left (Baseline):} off-diagonal entries below the diagonal drop
    sharply, quantifying forgetting.
    \textbf{Right (DCF-CL):} each column is constant from its training phase
    onward --- exact restore via atom swapping.}
  \label{fig:heatmap}
\end{figure}

\subsection{Memory Footprint}

Figure~\ref{fig:memory} compares the total stored parameters for both
approaches.  The baseline requires storing five independent models, totalling
\placeholder{base\_mem} parameters.  DCF-CL stores one shared set of
coefficients, decoder, and batch normalisation (\placeholder{shared}
parameters) plus five atom snapshots of 162 parameters each (810 total),
for a grand total of \placeholder{dcf\_mem} parameters --- a
\placeholder{ratio}$\times$ compression.

\begin{figure}[h]
  \centering
  \includegraphics[width=0.55\textwidth]{fig_p3v2_memory_footprint.pdf}
  \caption{Memory footprint for 5 tasks.
    Baseline stores 5 full models.
    DCF-CL stores one shared model plus 5 lightweight atom snapshots (162
    scalars each), achieving $\sim$\placeholder{ratio}$\times$ compression.}
  \label{fig:memory}
\end{figure}

%%=============================================================================
\section{Discussion}
\label{sec:discussion}
%%=============================================================================

\subsection{Why Storing Subspaces Prevents Catastrophic Forgetting}
\label{sec:disc_forget}

Catastrophic forgetting occurs because gradient updates for a new task
modify shared parameters in directions that are beneficial for the current
task but destructive for previously learned tasks.  The standard
continual-learning literature addresses this through regularisation
(EWC \cite{kirkpatrick2017overcoming}, SI), replay buffers, or parameter
isolation.  Our approach falls squarely in the \emph{parameter isolation}
category, and achieves the strongest possible guarantee: \textbf{exactly zero
forgetting}, not merely reduced forgetting.

The mechanism is algebraically simple.  Let
$\theta = (\Psi, \Phi)$ denote the full parameter vector, split into atoms
$\Psi = \{\psivec_k\}$ and everything else
$\Phi = \{\avec_k, \theta_\text{dec}, \theta_\text{BN}\}$.  After Task~0, we
fix $\Phi$ permanently.  For Task~$t$, we learn $\Psi^{(t)}$ while $\Phi$
is frozen, producing the task-specific model
$\theta^{(t)} = (\Psi^{(t)}, \Phi)$.  At inference time, restoring
$\Psi^{(j)}$ recovers $\theta^{(j)}$ \emph{exactly}, because $\Phi$ has not
changed.

The key insight is that there is \emph{no interference between task
subspaces}: the atoms for Task~$j$ live in their own copy of parameter memory
and are never modified by training on any other task.  Unlike
regularisation-based methods (which approximate ``importance'' and still permit
small drifts) or replay methods (which approximately re-learn old tasks), atom
swapping provides a \emph{hard guarantee} of zero forgetting --- the model for
Task~$j$ is bitwise identical to what it was immediately after Task~$j$'s
training.

This guarantee comes at a cost: the shared components $\Phi$ are frozen after
Task~0 and never adapted to later tasks.  If the reconstruction pipeline
learned on Task~0 is insufficiently general, later tasks may achieve lower
PSNR than they would with full end-to-end training.  Our experimental results
(Table~\ref{tab:final_psnr}) indicate that this cost is modest for the
Split-MNIST setting, where the shared blur kernel and similar image statistics
across digit classes make the Task-0 pipeline broadly applicable.

\subsection{Extension to Spatially Varying PSF Deblurring}
\label{sec:disc_svpsf}

The Split-MNIST experiment demonstrates the atom-swapping principle in a
controlled setting with discrete tasks.  My research on non-blind
deconvolution for spatially varying PSFs requires extending this principle to
a fundamentally different regime: \emph{continuous} blur conditions that
vary spatially within a single image.

\paragraph{From discrete tasks to continuous conditions.}
In our experiment, each task corresponds to a discrete set of digit classes,
and we store a separate atom snapshot $\Psi^{(t)}$ for each.  In SV-PSF
deblurring, a ``task'' corresponds to a specific local PSF condition
$h(\mathbf{x})$ at spatial position $\mathbf{x}$.  The PSF varies smoothly
across the image plane --- there are infinitely many ``tasks'' parameterised
by position and optical state.  Clearly, we cannot store infinitely many
discrete atom sets.

\paragraph{Filter atom swapping as inspiration.}
The Split-MNIST experiment validates that a powerful shared reconstruction
pipeline (coefficients + decoder) can be combined with lightweight
task-specific atoms to produce high-quality reconstructions.  This suggests
a design pattern for SV-PSF deblurring: train a \textbf{shared universal
deconvolution network and decoder} (analogous to the frozen $\Phi$), and
generate \textbf{spatially varying atoms} (analogous to $\Psi^{(t)}$) that
adapt the network to the local blur condition at each spatial position.

\paragraph{Connection to predictive interpolation.}
Since we operate in the \emph{non-blind} setting (the PSF is known or
estimated from calibration data), we have access to physical priors that
describe the blur at each position.  Rather than storing discrete atom
snapshots, we can train a lightweight \textbf{hypernetwork}
$f_\phi : h(\mathbf{x}) \mapsto \Psi(\mathbf{x})$ that predicts the
local atoms directly from the known local PSF:
\begin{equation}
  \psivec_k(\mathbf{x}) = f_\phi\bigl(h(\mathbf{x})\bigr)_k,
  \qquad k = 1, \ldots, K.
\end{equation}
This hypernetwork needs to predict only $K \times l^2$ values per spatial
position (e.g., $6 \times 9 = 54$ for our architecture), making it
extremely lightweight.  Because the PSF field $h(\mathbf{x})$ is spatially
smooth, the predicted atoms will also vary smoothly, and the hypernetwork
can be very small (a few fully connected layers) while still capturing the
relevant variation.

\paragraph{The vision: single-forward-pass SV-PSF deblurring.}
By feeding these generated, spatially varying atoms
$\Psi(\mathbf{x})$ into the shared deconvolution network, we can solve the
spatially varying deblurring problem \emph{in a single forward pass} with
an extremely low parameter cost.  The resulting architecture would:
\begin{enumerate}[leftmargin=*]
  \item Accept as input: a blurred image $y$ and a PSF field
        $\{h(\mathbf{x})\}$.
  \item Use the hypernetwork to predict spatially varying atoms
        $\Psi(\mathbf{x})$ from the PSF field.
  \item Inject these atoms into the shared DCF encoder, producing
        position-adapted feature maps.
  \item Decode the feature maps through the shared decoder to produce
        the restored image.
\end{enumerate}
This bypasses the need to train massive dynamic convolutions or
spatially varying filter banks, and instead leverages the low-dimensional
atom subspace structure to achieve spatial adaptivity at minimal cost.
The Split-MNIST experiment provides the proof of concept that the
shared-pipeline-plus-task-specific-atoms paradigm works; the hypernetwork
extension turns it into a practical architecture for real-world SV-PSF
deblurring.

%%=============================================================================
%% REFERENCES
%%=============================================================================

\begin{thebibliography}{9}

\bibitem{qiu2018dcfnet}
Q.~Qiu, X.~Cheng, and G.~Sapiro.
\newblock {DCFNet}: Deep neural network with decomposed convolutional filters.
\newblock In \emph{ICML}, 2018.

\bibitem{miao2021continual}
Z.~Miao, Z.~Wang, W.~Chen, and Q.~Qiu.
\newblock Continual learning with filter atom swapping.
\newblock In \emph{ICLR}, 2022.

\bibitem{french1999catastrophic}
R.~M. French.
\newblock Catastrophic forgetting in connectionist networks.
\newblock \emph{Trends in Cognitive Sciences}, 3(4):128--135, 1999.

\bibitem{kirkpatrick2017overcoming}
J.~Kirkpatrick et al.
\newblock Overcoming catastrophic forgetting in neural networks.
\newblock \emph{Proceedings of the National Academy of Sciences},
  114(13):3521--3526, 2017.

\bibitem{wang2020dictionary}
Z.~Wang, X.~Cheng, G.~Sapiro, and Q.~Qiu.
\newblock A dictionary approach to domain-invariant learning in deep networks.
\newblock In \emph{NeurIPS}, 2020.

\end{thebibliography}

\end{document}
